\subsection{定义与例子}

回顾（Ch. V, §6, Nos. 1 and 4；Ch. VI, §2, No.~10）：若 X 和 Y 是局部紧空间，\(\mu\) 是 X 上的测度，\(\varphi\) 是从 X 到 Y 的映射，则称 \(\varphi\) 是 \(\mu\)-适当的，如果：\(a)\) \(\varphi\) 是 \(\mu\)-可测的；\(b)\) 对 Y 的每个紧子集 K，\(\overline{\varphi}^1 (\mathbf{K})\) 都是本质上 \(\mu\)-可积的。于是 Y 上的像测度 \(\nu = \varphi (\mu)\) 存在，并具有如下性质：对于函数 \(\mathbf{f}\)，定义在 Y 上且取值于 Banach 空间或 \(\overline{\mathbf{R}}\)，它关于 \(\nu\) 本质可积，当且仅当 \(\mathbf{f} \circ \varphi\) 关于 \(\mu\) 本质可积；此时，

\[
\int_ {\mathrm{Y}} \mathbf {f} (y) d \nu (y) = \int_ {\mathrm{X}} \mathbf {f} (\varphi (x)) d \mu (x).
\]

定义 1. --- 设 \(X_1, \ldots, X_n\) 为局部紧空间，\(\mu_i\) 为 \(X_i\) 上的测度（\(1 \leqslant i \leqslant n\)）；令 \(X\) 为各 \(X_i\) 的乘积，\(\mu\) 为各 \(\mu_i\) 的乘积。令 \(\varphi\) 为从 \(X\) 到局部紧空间 \(Y\) 的映射。若序列 \((\mu_i)\) 是 \(\varphi\)-可卷积的，或者说 \(\mu_1, \ldots, \mu_n\) 是 \(\varphi\)-可卷积的，是指 \(\varphi\) 是 \(\mu\)-适当的；此时，称 \(\nu = \varphi(\mu)\) 为 \(\mu\) 在 \(\varphi\) 下的像，也称作 \(\mu_i\) 关于 \(\varphi\) 的卷积积，并记为 \(*^{\varphi(\mu_i)}_{1 \leqslant i \leqslant n}\)、\(\underset{i=1}{\overset{n}{*}} \mu_i\) 或 \(\mu_1 * \mu_2 * \cdots * \mu_n\)。

后两个记号当然只在 \(\varphi\) 不会引起歧义时使用。

设 f 为 Y 上取值于 Banach 空间或 \(\overline{\mathbf{R}}\) 的函数。f 关于 \(\mu_{1}*\cdots*\mu_{n}\) 本质可积，当且仅当函数

\[
\left(x _ {1}, \dots , x _ {n}\right) \mapsto \mathbf {f} \left(\varphi \left(x _ {1}, \dots , x _ {n}\right)\right)
\]

关于 \(\mu_{1}\otimes\mu_{2}\otimes\cdots\otimes\mu_{n}\) 本质可积；此时

\[
\int \mathbf {f} d (\mu_ {1} * \dots * \mu_ {n}) = \int \mathbf {f} \big (\varphi (x _ {1}, \ldots , x _ {n}) \big) d \mu_ {1} (x _ {1}) \ldots d \mu_ {n} (x _ {n}),\tag{1}
\]

这是一个可在取 \(\mu_{1}*\cdots *\mu_{n}\) 时视为其定义的公式，取 \(\mathbf{f}\in\mathcal{K}(\mathrm{Y})\) 时即如此。

由定义立即得到，\(\mu_{i}\) 可卷积，当且仅当 \(|\mu_{i}|\) 可卷积。在此情形下，

\[
\left| \varphi \left(\mu_ {1} \otimes \dots \otimes \mu_ {n}\right) \right| \leqslant \varphi \left(\left| \mu_ {1} \otimes \dots \otimes \mu_ {n} \right|\right) = \varphi \left(\left| \mu_ {1} \right| \otimes \dots \otimes \left| \mu_ {n} \right|\right)
\]

（Ch. VI, §2, No.~10），即

\[
\left| \begin{array}{c} * \\ i \end{array} \mu_ {i} \right| \leqslant * _ {i} | \mu_ {i} |.\tag{2}
\]

若 \(\mu_{i}\) 可卷积且为正测度，并且 \(\nu_{i}\) 是 \(X_{i}\) 上满足 \(0 \leqslant \nu_{i} \leqslant \mu_{i}\) 的测度，则 \(\nu_{i}\) 可卷积，并且

\[
\underset {i} {\ast} \nu_ {i} \leqslant \underset {i} {\ast} \mu_ {i}.
\]

设 \(\mu_1, \mu_2, \ldots, \mu_n\) 可卷积，且 \(\mu_1', \mu_2, \ldots, \mu_n\) 也可卷积（\(\mu_1'\) 是 \(X_1\) 上的测度）。根据 Ch. V, §6, No.~3, 命题 6 的推论 1，\(\mu_1 + \mu_1', \mu_2, \ldots, \mu_n\) 可卷积，并且

\[
\left(\mu_ {1} + \mu_ {1} ^ {\prime}\right) * \mu_ {2} * \dots * \mu_ {n} = \mu_ {1} * \mu_ {2} * \dots * \mu_ {n} + \mu_ {1} ^ {\prime} * \mu_ {2} * \dots * \mu_ {n}.
\]

例子。--- 1) 对任意 \(\varphi\)，测度 \(\varepsilon_{x_i}\)（其中 \(x_i \in X_i\) 对 \(1 \leqslant i \leqslant n\) 成立）总是可卷积的，其卷积积为 \(\varepsilon_y\)，其中 \(y = \varphi(x_1, x_2, \ldots, x_n)\)。因此，若每个 \(\mu_i\) 都有有限支集，则 \(\mu_i\) 可卷积，且 \(\mu_1 * \cdots * \mu_n\) 有有限支集。特别地，令 M 为赋有局部紧拓扑的幺半群 \(^1\)；若取 \(\varphi\) 为 M 中的复合律，则 M 上具有有限支集的测度在卷积下构成一个代数，这个代数恰好就是幺半群 M 的代数（视所考虑的是实测度还是复测度而分别取在 R 或 C 上）（A, III, §2, No.~6）。

\begin{enumerate}
\def\labelenumi{\arabic{enumi})}
\setcounter{enumi}{1}
\tightlist
\item
  设 M 为赋有离散拓扑的幺半群；假设对每个 \(m \in M\)，只有有限多个 \((m', m'') \in M \times M\) 满足 \(m'm'' = m\)；这等价于说 M 中的复合律是从 \(M \times M\) 到 M 的适当映射；于是 M 上的测度在卷积下构成一个代数，该代数正是幺半群 M 的全代数（A, III, §2, No.~10）；我们注意以下两个特殊情形：
\end{enumerate}

\begin{enumerate}
\def\labelenumi{\alph{enumi})}
\tightlist
\item
  M = N，复合律为加法。对 N 上的每个测度 \(\mu\)，令其对应于形式级数
\end{enumerate}

\[
\mathrm{S} (\mu) = \sum_ {n = 0} ^ {\infty} \mu (\{n \}) t ^ {n}
\]

其中 t 是未定元。则 \(S(\mu * \mu') = S(\mu)S(\mu')\)。对于任意多个未定元的形式级数，也有类似结论。

\(^{*}b)\) M = N \(^{*}\)，复合律为乘法。对每个测度 \(\mu\)（定义在 N \(^{*}\) 上），令其对应于形式 Dirichlet 级数

\[
\mathrm{D} (\mu) = \sum_ {n = 1} ^ {\infty} \mu (\{n \}) n ^ {- s}.
\]

则 \(\mathrm{D}(\mu * \mu') = \mathrm{D}(\mu) \mathrm{D}(\mu')\)。

\begin{enumerate}
\def\labelenumi{\arabic{enumi})}
\setcounter{enumi}{2}
\tightlist
\item
  设 X、Y、Z 为局部紧空间，\(\varphi\) 为从 \(X \times Y\) 到 Z 的连续映射。若 \(x \in X\)，且 \(\mu\) 是 Y 上的测度，则说 \(\varepsilon_{x}\) 与 \(\mu\) 关于 \(\varphi\) 可卷积，就是说从 Y 到 Z 的映射 \(\varphi(x,\cdot)\) 是 \(\mu\)-适当的。于是有 \(\varepsilon_{x} * \mu = \varphi(x,\cdot)(\mu)\)。
\end{enumerate}

